\documentclass[12pt]{article}
\usepackage[utf8]{inputenc}
\usepackage[T1]{fontenc}
\usepackage{hyperref}
\title{Guwa Language Resources for K-12 Curriculum Development}
\author{Compiled by Claude Sonnet 5.5 (claude-sonnet-5-5) for Andr\'as Kornai}
\date{}
\begin{document}
\maketitle

\section{Introduction}
Guwa (also spelt Goa or Koa; ISO 639-3 xgw, Glottocode guwa1242) is an extinct and very sparsely attested Aboriginal language of south-western Queensland. It is listed as Pama-Nyungan, unclassified or possibly Maric, and appears to have been close to Yanda. There are no known speakers, and the sources I found differ on the date of extinction and on the exact classification. Little exists beyond a linguistic salvage description and catalogue entries. A heritage curriculum would therefore have to be built with the descendant community and would rest on very limited material. I found no school programme, primer, audio or community language programme for Guwa.

\section{Primary Educational Resources}
I found no textbooks, primers or classroom materials for Guwa.

\section{Linguistic and Reference Materials}
\begin{itemize}
\item Blake, Barry and Gavan Breen (1990), ``Guwa'', pp.\ 108--144 in Gavan Breen (ed.), \emph{Salvage Studies of Western Queensland Aboriginal Languages}, Pacific Linguistics B-105. The volume is open access (CC BY-SA 4.0) from the ANU repository. In the search results I saw the PHOIBLE record of the chapter, \url{https://phoible.org/sources/309499}, and ANU repository links to the volume, \url{https://openresearch-repository.anu.edu.au/bitstreams/8c107f99-568b-4347-98ac-e823afe0ca3d/download}. I did not confirm which ANU file is the full volume and which is a part of it. I found no description of the chapter's contents beyond its title and page range.
\item Wikipedia, Guwa language: \url{https://en.wikipedia.org/wiki/Guwa_language}
\item Endangered Languages Project bibliography for Guwa, which lists 0 speakers: \url{https://endangeredlanguages.com/lang/6784/bibliography?hl=de}
\item ASJP database entry (word list data): \url{https://asjp.clld.org/languages/GUWAMU}
\end{itemize}

\section{Religious and Community Resources}
I found no scripture translations or community organisations specific to Guwa.

\section{Digital and Audio Resources}
I found no audio, video or software resources for Guwa.

\section{Curriculum Development Context}
Because the language is extinct and nearly unattested, any curriculum would be a heritage and awareness programme, not a speaker-based immersion course. The Blake and Breen chapter is the main documentary source and should be used only with the guidance of the relevant Aboriginal community, as the ANU repository carries a cultural advice notice. Related neighbouring languages such as Yanda and Guwar are separate and should not be conflated with Guwa.

\end{document}